\begin{textsample}{POS Dim 2 – vem\_ed\_32 – Score 7.00 – vem\_ed\_32/t172.txt}  \label{ex:f2_pos_046}
Intercâmbios Virtuais como Ponte de Empatia e Inovação no Ensino Artigo de Opinião Após ser convidada para participar de uma mesa-redonda no início do mês de outubro deste ano, \textbf{pedi} que alguns alunos gravassem vídeos relatando suas experiências em Projetos Colaborativos Internacionais (PCIs / CGESG), conhecidos mundialmente como Collaborative Online International Learning (COIL).
Como professora de inglês na Fatec Guaratinguetá, já participei de PCIs com instituições dos Estados Unidos, Armênia, Filipinas, México e Israel.
Essas parcerias sempre me pareceram uma oportunidade pedagógica, mas ouvir os estudantes falando me fez \textbf{perceber} o impacto humano e transformador desses projetos.
Muitos alunos relataram terem se sentido mais confiantes depois de interagir com colegas de outros países.
Alguns disseram que, graças ao PCI, se prepararam para \textbf{entrevistas} de emprego em inglês e conseguiram oportunidades profissio- nais.
Outros contaram que o intercâmbio virtual reacendeu seu interesse pela língua inglesa: não mais apenas um conteúdo da escola, mas algo com significado real em suas vidas.
Um dos depoimentos mais marcantes foi de um aluno que lutava contra a depressão.
Ele disse que, depois de participar de um PCI, redescobriu sua paixão pelo inglês, criou um canal no YouTube e encontrou uma nova razão para continuar estudando.
Para mim, esse tipo de transformação mostra que o PCI / COIL é mais do que ensino: é uma força de mudança social.
No nível docente, também experimento um profundo crescimento.
Planejar cursos com parceiros internacionais exige empatia, cooperação e adaptação — e nos deixa mais abertos às realidades dos outros.
Mesmo com a distância física, percebo que muitos desafios educacionais que enfrentamos (burocracia, limita- ção tecnológica, engajamento) são compartilhados por professores ao redor do mundo.
Isso fortalece a sensação de comunidade profissional global.
Além disso, os projetos ajudam a cultivar cidadania global.
Os alunos \textbf{percebem} que suas vozes têm relevância além de sua região, que podem dialogar com outras culturas e refletir sobre suas próprias identidades.
Esses intercâmbios os incentivam a pensar globalmente e \textbf{agir} localmente.
Re- conheço que há obstáculos reais: fusos horários, infraestrutura desigual e restrições tecnológicas.
Mas acredito que esses desafios se transformam em oportunidades.
Ao buscar ajustes, somos obrigados a repensar métodos, cronogramas e abordagens, e isso gera \textbf{inovação} e resiliência.
Por fim, afirmo que os projetos COIL merecem apoio institucional mais robusto: é \textbf{preciso} tempo para planejamento, recursos tecnológicos, capacitação docente.
Esse investimento não é apenas estratégico: é humanitário.
Quando estudantes relatam que o COIL lhes deu confiança, propósito ou alegria, \textbf{percebemos} que a internacionalização virtual não serve apenas para “fazer intercâmbio” — ela pode mudar vidas.
E, para nós, professores, é a confirmação de que ensinar é, antes de tudo, ligar pontes humanas.
Referências HACKETT, S.; DAWSON, M.; JANSSEN, J.; VAN TARTWIJK, J.
Defining Collaborative Online International Learning (COIL) and distinguishing it from Virtual Exchange.
TechTrends, v. 68, n. 6, p. 1078-1094, 2024.
Disponível em: https://doi.org/10.1007/s11528-024-01000-w.
Acesso em: 28 set. 2025.
HACKETT, S.; JANSSEN, J.; BEACH, P.; PERREAULT, M.; BEELEN, J.; VAN TARTWIJK, J.
The effectiveness of Collaborative Online International Learning (COIL) on intercultural competence development in higher education.
International Journal of Educational Technology in Higher Education, v. 20, n. 1, art. 5, 2023.
Disponível em: https://doi.org/10.1186/s41239-022-00373-3.
Acesso em: 28 set. 2025.
RAJAGOPALAN, Kanavillil.
Vencer barreiras e emergir das adversidades com pleno êxito, sempre com o pé no chão.
In: LIMA, D.
C. de (Org.).
Inglês em escolas públicas não funciona?
Uma questão, \textbf{múltiplos} olhares.
São Paulo: Parábola, 2011, p. 55–65.

% matched lemmas: agir, entrevista, inovação, múltiplo, pedir, perceber, preciso
\end{textsample}
